% problem_id: S001-proposition-poincare-duality
% source_id: S001 (Stacks Project)
% source_file: derham.tex
% statement_locator: derham.tex lines 4472--4488
% proof_locator: derham.tex lines 4490--4662
% env: proposition
% status: text-extracted-verbatim (difficulty judgement pending, written by hand)
%
% The statement and proof below are the author's own text, extracted verbatim.
% Nothing is paraphrased, shortened, or supplemented.

\begin{proposition}
\label{proposition-poincare-duality}
Let $k$ be a field. Let $X$ be a nonempty smooth proper scheme over $k$
equidimensional of dimension $d$. There exists a $k$-linear map
$$
t : H^{2d}_{dR}(X/k) \longrightarrow k
$$
unique up to precomposing by multiplication by a unit of
$H^0(X, \mathcal{O}_X)$ with the following property: for all $i$ the pairing
$$
H^i_{dR}(X/k) \times H_{dR}^{2d - i}(X/k)
\longrightarrow
k, \quad
(\xi, \xi') \longmapsto t(\xi \cup \xi')
$$
is perfect.
\end{proposition}

\begin{proof}
By the Hodge-to-de Rham spectral sequence
(Section \ref{section-hodge-to-de-rham}), the vanishing
of $\Omega^i_{X/k}$ for $i > d$, the vanishing in
Cohomology, Proposition \ref{cohomology-proposition-vanishing-Noetherian}
and the results of Lemmas \ref{lemma-bottom-part-degenerates} and
\ref{lemma-top-part-degenerates}
we see that $H^0_{dR}(X/k) = H^0(X, \mathcal{O}_X)$
and $H^d(X, \Omega^d_{X/k}) = H_{dR}^{2d}(X/k)$.
More precisely, these identifications come from the maps
of complexes
$$
\Omega^\bullet_{X/k} \to \mathcal{O}_X[0]
\quad\text{and}\quad
\Omega^d_{X/k}[-d] \to \Omega^\bullet_{X/k}
$$
Let us choose $t : H_{dR}^{2d}(X/k) \to k$ which via this identification
corresponds to a $t$ as in Lemma \ref{lemma-duality-hodge}.
Then in any case we see that the pairing displayed in the lemma
is perfect for $i = 0$.

\medskip\noindent
Denote $\underline{k}$ the constant sheaf with value $k$ on $X$.
Let us abbreviate $\Omega^\bullet = \Omega^\bullet_{X/k}$.
Consider the map (\ref{equation-wedge}) which in our situation reads
$$
\wedge :
\text{Tot}(\Omega^\bullet \otimes_{\underline{k}} \Omega^\bullet)
\longrightarrow
\Omega^\bullet
$$
For every integer $p = 0, 1, \ldots, d$ this map
annihilates the subcomplex
$\text{Tot}(\sigma_{> p} \Omega^\bullet \otimes_{\underline{k}}
\sigma_{\geq d - p} \Omega^\bullet)$ for degree reasons.
Hence we find that the restriction of $\wedge$ to the subcomplex
$\text{Tot}(\Omega^\bullet \otimes_{\underline{k}}
\sigma_{\geq d - p}\Omega^\bullet)$ factors through a map of complexes
$$
\gamma_p :
\text{Tot}(\sigma_{\leq p} \Omega^\bullet \otimes_{\underline{k}}
\sigma_{\geq d - p} \Omega^\bullet)
\longrightarrow
\Omega^\bullet
$$
Using the same procedure as in Section \ref{section-cup-product} we obtain
cup products
$$
H^i(X, \sigma_{\leq p} \Omega^\bullet) \times
H^{2d - i}(X, \sigma_{\geq d - p}\Omega^\bullet)
\longrightarrow
H_{dR}^{2d}(X, \Omega^\bullet)
$$
We will prove by induction on $p$ that these cup products via $t$
induce perfect pairings between $H^i(X, \sigma_{\leq p} \Omega^\bullet)$
and $H^{2d - i}(X, \sigma_{\geq d - p}\Omega^\bullet)$. For $p = d$
this is the assertion of the proposition.

\medskip\noindent
The base case is $p = 0$. In this case we simply obtain the pairing
between $H^i(X, \mathcal{O}_X)$ and $H^{d - i}(X, \Omega^d)$ of
Lemma \ref{lemma-duality-hodge} and the result is true.

\medskip\noindent
Induction step. Say we know the result is true for $p$. Then
we consider the distinguished triangle
$$
\Omega^{p + 1}[-p - 1] \to
\sigma_{\leq p + 1}\Omega^\bullet \to
\sigma_{\leq p}\Omega^\bullet \to
\Omega^{p + 1}[-p]
$$
and the distinguished triangle
$$
\sigma_{\geq d - p}\Omega^\bullet \to
\sigma_{\geq d - p - 1}\Omega^\bullet \to
\Omega^{d - p - 1}[-d + p + 1] \to
(\sigma_{\geq d - p}\Omega^\bullet)[1]
$$
Observe that both are distinguished triangles in the homotopy category
of complexes of sheaves of $\underline{k}$-modules; in particular the
maps $\sigma_{\leq p}\Omega^\bullet \to \Omega^{p + 1}[-p]$ and
$\Omega^{d - p - 1}[-d + p + 1] \to (\sigma_{\geq d - p}\Omega^\bullet)[1]$
are given by actual maps of complexes, namely using the differential
$\Omega^p \to \Omega^{p + 1}$ and the differential
$\Omega^{d - p - 1} \to \Omega^{d - p}$.
Consider the long exact cohomology sequences associated to these
distinguished triangles
$$
\xymatrix{
H^{i - 1}(X, \sigma_{\leq p}\Omega^\bullet) \ar[d]_a \\
H^i(X, \Omega^{p + 1}[-p - 1]) \ar[d]_b \\
H^i(X, \sigma_{\leq p + 1}\Omega^\bullet) \ar[d]_c \\
H^i(X, \sigma_{\leq p}\Omega^\bullet) \ar[d]_d \\
H^{i + 1}(X, \Omega^{p + 1}[-p - 1])
}
\quad\quad
\xymatrix{
H^{2d - i  + 1}(X, \sigma_{\geq d - p}\Omega^\bullet) \\
H^{2d - i}(X, \Omega^{d - p - 1}[-d + p + 1]) \ar[u]_{a'} \\
H^{2d - i}(X, \sigma_{\geq d - p - 1}\Omega^\bullet) \ar[u]_{b'} \\
H^{2d - i}(X, \sigma_{\geq d - p}\Omega^\bullet) \ar[u]_{c'} \\
H^{2d - i - 1}(X, \Omega^{d - p - 1}[-d + p + 1]) \ar[u]_{d'}
}
$$
By induction and Lemma \ref{lemma-duality-hodge}
we know that the pairings constructed above between the
$k$-vectorspaces on the first, second, fourth, and fifth
rows are perfect. By the $5$-lemma, in order to show that
the pairing between the cohomology groups in the middle row
is perfect, it suffices to show that the pairs
$(a, a')$, $(b, b')$, $(c, c')$, and $(d, d')$
are compatible with the given pairings (see below).

\medskip\noindent
Let us prove this for the pair $(c, c')$. Here we observe simply
that we have a commutative diagram
$$
\xymatrix{
\text{Tot}(\sigma_{\leq p} \Omega^\bullet \otimes_{\underline{k}}
\sigma_{\geq d - p} \Omega^\bullet) \ar[d]_{\gamma_p} &
\text{Tot}(\sigma_{\leq p + 1} \Omega^\bullet \otimes_{\underline{k}}
\sigma_{\geq d - p} \Omega^\bullet) \ar[l] \ar[d] \\
\Omega^\bullet &
\text{Tot}(\sigma_{\leq p + 1} \Omega^\bullet \otimes_{\underline{k}}
\sigma_{\geq d - p - 1} \Omega^\bullet) \ar[l]_-{\gamma_{p + 1}}
}
$$
Hence if we have $\alpha \in H^i(X, \sigma_{\leq p + 1}\Omega^\bullet)$
and $\beta \in H^{2d - i}(X, \sigma_{\geq d - p}\Omega^\bullet)$
then we get
$\gamma_p(\alpha \cup c'(\beta)) = \gamma_{p + 1}(c(\alpha) \cup \beta)$
by functoriality of the cup product.

\medskip\noindent
Similarly for the pair $(b, b')$ we use the commutative diagram
$$
\xymatrix{
\text{Tot}(\sigma_{\leq p + 1} \Omega^\bullet \otimes_{\underline{k}}
\sigma_{\geq d - p - 1} \Omega^\bullet) \ar[d]_{\gamma_{p + 1}} &
\text{Tot}(\Omega^{p + 1}[-p - 1] \otimes_{\underline{k}}
\sigma_{\geq d - p - 1} \Omega^\bullet) \ar[l] \ar[d] \\
\Omega^\bullet &
\Omega^{p + 1}[-p - 1]
\otimes_{\underline{k}}
\Omega^{d - p - 1}[-d + p + 1] \ar[l]_-\wedge
}
$$
and argue in the same manner.

\medskip\noindent
For the pair $(d, d')$ we use the commutative diagram
$$
\xymatrix{
\Omega^{p + 1}[-p] \otimes_{\underline{k}}
\Omega^{d - p - 1}[-d + p] \ar[d] &
\text{Tot}(\sigma_{\leq p}\Omega^\bullet \otimes_{\underline{k}}
\Omega^{d - p - 1}[-d + p]) \ar[l] \ar[d] \\
\Omega^\bullet &
\text{Tot}(\sigma_{\leq p}\Omega^\bullet \otimes_{\underline{k}}
\sigma_{\geq d - p}\Omega^\bullet) \ar[l]
}
$$
and we look at cohomology classes in
$H^i(X, \sigma_{\leq p}\Omega^\bullet)$ and
$H^{2d - i}(X, \Omega^{d - p - 1}[-d + p])$.
Changing $i$ to $i - 1$ we get the result for the pair $(a, a')$
thereby finishing the proof that our pairings are perfect.

\medskip\noindent
We omit the argument showing the uniqueness of $t$ up to
precomposing by multiplication by a unit in $H^0(X, \mathcal{O}_X)$.
\end{proof}
